\Needspace{8\baselineskip}
\Article{201}
Zum Problem, das wir in den vorhergehenden Artikeln behandelt haben, fügen wir noch folgende Bemerkungen hinzu.

1) Wir haben gezeigt, wie man die Gleichung \(tt-Duu=mm\) in allen Fällen löst, in denen \(m\) der größte gemeinsame Teiler dreier Zahlen \(M,2N,P\) mit \(NN-MP=D\) ist. Es lohnt sich, alle Zahlen zu bestimmen, die solche Teiler sein können, also alle Werte von \(m\) für einen gegebenen Wert von \(D\). Wir setzen \(D=nnD'\), sodass \(D'\) völlig frei von quadratischen Faktoren ist. Das erreichen wir, indem wir für \(nn\) das größte Quadrat nehmen, das \(D\) teilt; hat \(D\) selbst schon keinen quadratischen Faktor, so müssen wir \(n=1\) setzen. Dann behaupte ich Folgendes.

\textit{Erstens}: Hat \(D'\) die Form \(4k+1\), so ist jeder Teiler von \(2n\) ein Wert von \(m\) und umgekehrt. Teilt nämlich \(g\) die Zahl \(2n\), so haben wir die Form \(\left(g,n,\frac{nn(1-D')}{g}\right)\) mit der Determinante \(D\), bei der die Zahlen \(g,2n,\frac{nn(D'-1)}{g}\) offenbar den größten gemeinsamen Teiler \(g\) haben (denn \(\frac{nn(D'-1)}{gg}=\frac{4nn}{gg}\cdot\frac{D'-1}{4}\) ist offenbar ganzzahlig). Nehmen wir umgekehrt an, \(g\) sei ein Wert von \(m\), nämlich der größte gemeinsame Teiler von \(M,2N,P\), und es gelte \(NN-MP=D\). Dann ist offenbar \(4D\) oder \(4nnD'\) durch \(gg\) teilbar. Daraus folgt, dass \(2n\) durch \(g\) teilbar sein muss. Würde nämlich \(g\) die Zahl \(2n\) nicht teilen, so hätten \(g\) und \(2n\) einen größten gemeinsamen Teiler kleiner als \(g\). Setzen wir diesen \(=\delta\) und
\SourceLeaf{198}
\(2n=\delta n',\ g=\delta g'\), so wäre \(n'n'D'\) durch \(g'g'\) teilbar, \(n'\) zu \(g'\) und daher auch \(n'n'\) zu \(g'g'\) teilerfremd; somit wäre auch \(D'\) durch \(g'g'\) teilbar, im Widerspruch zur Voraussetzung, dass \(D'\) von jedem quadratischen Faktor frei ist.

\textit{Zweitens}: Hat \(D'\) die Form \(4k+2\) oder \(4k+3\), so ist jeder Teiler von \(n\) ein Wert von \(m\), und umgekehrt teilt jeder Wert von \(m\) die Zahl \(n\). Teilt nämlich \(g\) die Zahl \(n\), so haben wir die Form \(\left(g,0,-\frac{nnD'}{g}\right)\) mit der Determinante \(=D\), bei der die Zahlen \(g,0,\frac{nnD'}{g}\) offenbar den größten gemeinsamen Teiler \(g\) haben. --- Nehmen wir dagegen an, \(g\) sei ein Wert von \(m\), nämlich der größte gemeinsame Teiler von \(M,2N,P\), und es gelte \(NN-MP=D\), so teilt wie oben \(g\) die Zahl \(2n\), oder \(\frac{2n}{g}\) ist ganzzahlig. Wäre dieser Quotient ungerade, so wäre das Quadrat \(\frac{4nn}{gg}\) \(\equiv1\ (\mathrm{mod.}\ 4)\) und daher \(\frac{4nnD'}{gg}\) entweder \(\equiv2\) oder \(\equiv3\ (\mathrm{mod.}\ 4)\). Nun gilt aber \(\frac{4nnD'}{gg}=\frac{4D}{gg}=\frac{4NN}{gg}-\frac{4MP}{gg}\equiv\frac{4NN}{gg}(\mathrm{mod.}\ 4)\), also wäre \(\frac{4NN}{gg}\) entweder \(\equiv2\) oder \(\equiv3\ (\mathrm{mod.}\ 4)\), \textit{ein Widerspruch}, da jedes Quadrat modulo \(4\) entweder zu null oder zu eins kongruent sein muss. Daher ist der Quotient \(\frac{2n}{g}\) notwendig gerade und somit \(\frac{n}{g}\) ganzzahlig, also \(g\) ein Teiler von \(n\).

Also ist \(1\) stets ein Wert von \(m\), oder die Gleichung \(tt-Duu=1\) lässt sich nach dem Vorhergehenden für jeden positiven nichtquadratischen Wert von \(D\) lösen; ferner ist \(2\) genau dann ein Wert von \(m\), wenn \(D\) entweder die Form \(4k\) oder \(4k+1\) hat.

2) Ist \(m\) größer als \(2\), aber eine zulässige Zahl, so können wir die Lösung von \(tt-Duu=mm\) auf die einer ähnlichen Gleichung zurückführen, in der \(m\) entweder \(1\) oder \(2\) ist. Setzen wir wie zuvor \(D=nnD'\), so teilt, wenn \(m\) die Zahl \(n\) teilt, auch \(mm\) die Zahl \(D\). Sind die kleinsten Werte von \(p,q\) in \(pp-\frac{D}{mm}qq=1\) gleich \(p=P,\ q=Q\), so sind die kleinsten Werte von \(t,u\) in \(tt-Duu=mm\) gleich \(t=mP,\ u=Q\). --- Teilt aber \(m\) die Zahl \(n\) nicht, so teilt es wenigstens \(2n\) und ist sicher gerade; \(\frac{4D}{mm}\) ist dann ganzzahlig. Finden wir in diesem Fall als kleinste Werte von \(p,q\) in \(pp-\frac{4D}{mm}qq=4\) die Werte \(p=P,\ q=Q\), so sind die kleinsten Werte von \(t,u\) in \(tt-Duu=mm\) gleich \(t=\frac{m}{2}P,\ u=Q\). --- In beiden Fällen lassen sich nach dieser Methode offenbar nicht nur aus den kleinsten Werten von \(p,q\) die kleinsten Werte von \(t,u\), sondern aus \textit{allen} Werten der ersten \textit{alle} Werte der zweiten ableiten.

3) Bezeichnen \(t^0,u^0;\ t',u';\ t'',u''\) usw. sämtliche positiven Werte von
\SourceLeaf{199}
\(t,u\) in \(tt-Duu=mm\) (wie im vorhergehenden Artikel), und sind gewisse Werte dieser Folge ihren ersten Werten modulo einer beliebigen gegebenen Zahl \(r\) kongruent, etwa \(t^\rho\equiv t\) (oder \(\equiv m\)), \(u^\rho\equiv u^0\) oder \(\equiv0\ (\mathrm{mod.}\ r)\), und zugleich die unmittelbar folgenden Werte den zweiten, etwa
\[
t^{\rho+1}\equiv t',\quad u^{\rho+1}\equiv u'\quad(\mathrm{mod.}\ r)
\]
so gilt auch
\[
t^{\rho+2}\equiv t'',\quad u^{\rho+2}\equiv u'';\quad t^{\rho+3}\equiv t''',\quad u^{\rho+3}\equiv u'''\text{ usw.}
\]
Das folgt leicht daraus, dass beide Folgen \(t^0,t',t''\) usw., \(u^0,u',u''\) usw. Rekursionen erfüllen: Denn wegen
\[
t''=\frac{2T}{m}t'-t^0,\quad t^{\rho+2}=\frac{2T}{m}t^{\rho+1}-t^\rho
\]
gilt
\[
t''\equiv t^{\rho+2}
\]
und ebenso bei den übrigen. --- Daraus folgt allgemein
\[
t^{h+\rho}\equiv t^h,\quad u^{h+\rho}\equiv u^h(\mathrm{mod.}\ r)
\]
wobei \(h\) eine beliebige Zahl bezeichnet, und noch allgemeiner, wenn
\[
\mu\equiv\nu\ (\mathrm{mod.}\ \rho),\quad\text{so gilt }t^\mu\equiv t^\nu,\quad u^\mu\equiv u^\nu(\mathrm{mod.}\ r).
\]

4) Die Bedingungen der vorhergehenden Bemerkung lassen sich aber immer erfüllen: Wir können stets einen Index \(\rho\) finden (für jeden gegebenen Modul \(r\)), für den gilt
\[
t^\rho\equiv t^0,\quad t^{\rho+1}\equiv t',\quad u^\rho\equiv u^0,\quad u^{\rho+1}\equiv u'
\]
Um das zu beweisen, bemerken wir Folgendes.

\textit{Erstens} lässt sich die dritte Bedingung immer erfüllen. Nach den Kriterien in (1) ist offenbar auch die Gleichung \(pp-rrDqq=mm\) lösbar. Sind die kleinsten positiven Werte von \(p,q\) (außer \(m,0\)) gleich \(P,Q\), so gehören zu den Werten von \(t,u\) offenbar auch \(t=P,\ u=rQ\). Daher kommen \(P,rQ\) in den Folgen \(t^0,t'\) usw., \(u^0,u'\) usw. vor; und wenn \(P=t^\lambda,\ rQ=u^\lambda\), so gilt \(u^\lambda\equiv0\equiv u^0(\mathrm{mod.}\ r)\). Außerdem sieht man leicht, dass zwischen \(u^0\) und \(u^\lambda\) kein Glied zu \(u^0\) modulo \(r\) kongruent ist.

\textit{Zweitens} ist klar: Sind außerdem die drei übrigen Bedingungen erfüllt, gilt also auch \(u^{\lambda+1}\equiv u',\ t^\lambda\equiv t^0,\ t^{\lambda+1}\equiv t'\), so brauchen wir nur \(\rho=\lambda\) zu setzen. Wenn
\SourceLeaf{200}
dagegen die eine oder andere dieser Bedingungen nicht gilt, so können wir, behaupte ich, sicher \(\rho=2\lambda\) setzen. Aus Gleichung~\([1]\) und den allgemeinen Formeln für \(t^e,u^e\) im vorhergehenden Artikel folgt nämlich
\[
t^{2\lambda}=\frac1m(t^\lambda t^\lambda+Du^\lambda u^\lambda)=\frac1m(mm+2Du^\lambda u^\lambda)
\]
und daher
\[
\frac{t^{2\lambda}-t^0}{r}=\frac{2Du^\lambda u^\lambda}{mr}
\]
Diese Größe ist ganzzahlig, weil nach Voraussetzung \(r\) die Zahl \(u^\lambda\) teilt, ebenso \(mm\) die Zahl \(4D\) und somit erst recht \(m\) die Zahl \(2D\). --- Ferner gilt \(u^{2\lambda}=\frac2m t^\lambda u^\lambda\), und wegen
\[
4t^\lambda t^\lambda=4Du^\lambda u^\lambda+4mm
\]
und damit der Teilbarkeit durch \(mm\) ist \(2t^\lambda\) durch \(m\) teilbar und somit \(u^{2\lambda}\) durch \(r\), also
\[
u^{2\lambda}\equiv u^0\ (\mathrm{mod.}\ r)
\]
Drittens finden wir
\[
t^{2\lambda+1}=t'+\frac{2Du^\lambda u^{\lambda+1}}{m}
\]
und da \(\frac{2Du^\lambda}{mr}\) aus demselben Grund ganzzahlig ist, gilt
\[
t^{2\lambda+1}\equiv t'\ (\mathrm{mod.}\ r)
\]
Schließlich finden wir
\[
u^{2\lambda+1}=u'+\frac{2t^{\lambda+1}u^\lambda}{m}
\]
und weil \(2t^{\lambda+1}\) durch \(m\) und \(u^\lambda\) durch \(r\) teilbar ist, gilt
\[
u^{2\lambda+1}\equiv u'\ (\mathrm{mod.}\ r).\quad \textit{Damit ist die Behauptung bewiesen.}
\]

Wozu die letzten beiden Bemerkungen dienen, wird sich im Folgenden zeigen.

\Needspace{30\baselineskip}
\Article{202}
Der besondere Fall des Problems, nämlich die Gleichung \(tt-Duu=1\) zu lösen, wurde schon von den Mathematikern des vorhergehenden Jahrhunderts behandelt. Der sehr scharfsinnige Fermat legte dieses Problem den englischen Analysten vor. Wallis nennt Brounker als Urheber der Lösung, die er in \textit{Alg. Cap.}~98, \textit{Opp. T.}~II \textit{p.}~418 ff. angibt; Ozanam nennt Fermat; und schließlich der bedeutende Euler, der das Problem in \textit{Comm. Petr.}~VI \textit{p.}~175, \textit{Comm. nov.}~XI \textit{p.}~28\HistoricalNote{*)}{In dieser Abhandlung wird der Algorithmus, den wir in Artikel~27 dargestellt haben, durch ähnliche Zeichen ausgedrückt; wir haben versäumt, das dort zu erwähnen.}, \textit{Algebra P.}~II \textit{p.}~226, \textit{Opusc. An.}~I \textit{p.}~310 behandelt, nennt Pell. Daher
\SourceLeaf{201}
haben einige Autoren das Problem als \textit{Pellsches Problem} bezeichnet. Alle diese Lösungen stimmen im Wesentlichen mit der überein, die wir erhalten, wenn wir in Artikel~198 die reduzierte Form mit \(a=1\) wählen. Dass aber das von ihnen vorgeschriebene Verfahren notwendig einmal \textit{endet}, also das Problem immer \textit{tatsächlich lösbar} ist, hat vor dem bedeutenden La Grange niemand streng\HistoricalNote{*)}{Was Wallis dazu an der genannten Stelle auf S.~427, 428 vorbringt, hat kein Gewicht. Sein Fehlschluss besteht darin, dass er auf S.~428, Zeile~4 voraussetzt, man könne zu einer gegebenen Größe \(p\) ganze Zahlen \(a,z\) finden, sodass \(\frac{z}{a}\) kleiner als \(p\) und der Fehlbetrag kleiner als der \textit{vorgegebene} ist. Das ist durchaus richtig, wenn der \textit{vorgegebene} Fehlbetrag eine \textit{gegebene Größe} ist; nicht aber, wenn er von \(a\) und \(z\) abhängt und somit veränderlich ist, wie im vorliegenden Fall.} bewiesen: \textit{Mélanges de la Soc. de Turin T.}~IV \textit{p.}~19 und eleganter in \textit{Hist. de l'Ac. de Berlin}, 1767, \textit{p.}~237. Diese Untersuchung steht auch in den schon mehrfach angeführten \textit{Ergänzungen zu Eulers Algebra}. Unsere Methode dagegen (die von völlig anderen Grundsätzen ausgeht und nicht auf \(m=1\) beschränkt ist) bietet meist mehrere Wege zur Lösung, da wir in Artikel~198 von jeder anderen reduzierten Form \((a,b,-a')\) ausgehen können.
